\part{Usage of the \texttt{tumbeamer} class}
\section{Basic settings}
\subsection{Class options}
\begin{frame}[fragile]{Class options}
  The class is loaded similar to any other class with the command
  \lstinline|\documentclass[<options>]{tumbeamer}|. Available options
  are key=value pairs and are summarized in the following table.

  \vspace{3mm}
  \begin{tabular}{lp{2cm}ll}
    \toprule
    option      & values             & default & Description               \\
    \midrule
    aspectratio & 169, 43            & 169     & choose the slide format   \\
    handout     & simple, 2on1, 4on1 & --      & enable handout mode with
    1, 2 or 4 slides per page                                              \\
    partpage    & true, false        & false   & enable slide at beginning
    of every part                                                          \\
    sectionpage & true, false        & true    & enable slide at beginning
    of every section                                                       \\
    \bottomrule
  \end{tabular}
  \vspace{.1mm}

  Additionally you can pass options to the base class \texttt{beamer}.
  Please refer to the \href{https://ctan.org/pkg/beamer}{beamer manual}
  for their usage.
\end{frame}

\subsection{Style setup}
\begin{frame}[fragile]{Style setup}
  For setting up the style of your presentation you use the command
  \lstinline|\TUMbeamersetup{<options>}|. The default setup is shown
  in the following listing.

  \begin{lstlisting}
    \TUMbeamersetup{
      title page   = TUM tower,
      part page    = TUM toc,
      section page = TUM toc,
      content page = TUM default,
      headline     = TUM threeliner,
      headline on  = {title page},
      footline     = TUM default,
      footline on  = {every page, title page=false},
    }
  \end{lstlisting}
\end{frame}

\begin{frame}{Style setup}
  \texttt{tumbeamer} uses a concept of \textit{framestyles}, which are
  collections of beamertemplate definitions combined to get a certain style.
  Initially we defined 4 different framestyles for \texttt{title page},
  \texttt{part page}, \texttt{section page}, and \texttt{content page},
  respectively. \\
  The predefined options for every framestyle are summarized on a following
  slide. \\\vspace{3mm}

  With the keys \texttt{headline} and \texttt{footline} you can choose from
  predefined styles.

  Enabling or disabling of headlines and footlines can be done on a per
  framestyle basis with the keys \texttt{headline on} and
  \texttt{footline on}. As value you can pass a list of framestyles or the
  alias \texttt{every page}, which enables the headline or footline on all 4
  framestyles.
\end{frame}

\begin{frame}{Predefined framestyles}
  List of predefined styles for \texttt{title page}, \texttt{part page},
  and \texttt{section page}.

  \vspace{3mm}
  \begin{tabular}{ll}
    \toprule
    possible values                                   & background color \\
    \midrule
    TUM default, TUM centered, TUM tower              & white            \\
    TUM blue, TUM blue centered, TUM blue tower       & blue             \\
    TUM shaded, TUM shaded centered, TUM shaded tower & blue shaded      \\
    TUM flags                                         & flag picture     \\
    \bottomrule
  \end{tabular}
  \vspace{3mm}

  For the section and part pages there is the additional style
  "TUM toc" available, which prints a table of content of the current
  part or section, respectively.
\end{frame}

\begin{frame}{Predefined framestyles}
  For content pages, there are fewer options possible (and reasonable).

  \vspace{3mm}
  \begin{tabular}{ll}
    TUM default    & white, frametitle $1\times$ logoheight below TUM logo \\
    TUM more space & white, frametitle at same height with TUM logo        \\
    TUM blue       & default layout with blue background                   \\
    TUM shaded     & default layout with shaded background                 \\
  \end{tabular}
\end{frame}

\begin{frame}[fragile]{Predefined styles for headlines and footlines}
  Possible options for headlines.

  \vspace{1mm}
  \begin{tabular}{ll}
    TUM empty      & no headline at all    \\
    TUM oneliner   & the TUM CD oneliner   \\
    TUM twoliner   & the TUM CD twoliner   \\
    TUM threeliner & the TUM CD threeliner \\
  \end{tabular}
  \vspace{3mm}

  Possible options for footlines. For the "TUM default" option, the user has
  to provide a custom footline with the command
  \lstinline|\footline{my footline text}|.

  \vspace{1mm}
  \begin{tabular}{ll}
    TUM empty    & no footline at all                                \\
    TUM default  & user defined footline and additional slide number \\
    TUM infoline & predefined footline with autor, shorttitle, date,
    and slide number                                                 \\
  \end{tabular}
\end{frame}

\section{Advanced settings}
\subsection{Handling custom framestyles}
\begin{frame}[fragile]{Handling custom framestyles}
  As explained above we use a concept of \texttt{framestyles} and we have
  predefined 4 framestyles which should cover most needs.
  If a user wants to have complete control over the style of a particular
  frame, we offer macros to define and load custom framestyles.

  \vspace{3mm}
  \lstinline|\defTUMframestyle{<name>}{<code>}| \\
  \hspace{5mm} Define a new framestyle <name> with the given <code>\\
  \lstinline|\setTUMframestyle{<name>}{<code>}| \\
  \hspace{5mm} Overwrite an existing framestyle <name> with the given <code>\\
  \lstinline|\addtoTUMframestyle{<name>}{<code>}| \\
  \hspace{5mm} Add the given <code> to an existing framestyle <name>\\
  \lstinline|\TUMframestyle{<name>}| \\
  \hspace{5mm} Load the framestyle <name> \\
\end{frame}

\begin{frame}[fragile]{Handling custom framestyles}
  The <code> that should go in these commands is intendet to load or define
  beamer templates. Please refer to the
  \href{https://ctan.org/pkg/beamer}{beamer manual} if you are not sure how
  to use these templates.

  To load the framestyle use the command before the frame
  \begin{lstlisting}
    .\TUMframestyle{section page}
    .\begin{frame}{my frametitle}
    .  my content
    .\end{frame}
  \end{lstlisting}
\end{frame}

\subsection{Color themes}
\begin{frame}[fragile]{Color themes}
  For color management we use the beamer framework to set and use colors.
  Please refer to the \href{https://ctan.org/pkg/beamer}{beamer manual}
  if you want to know more about setting and using colors.

  But if you use the standard \lstinline|\setbeamercolor{<name>}{<color>}|
  command offered by beamer in the preamble of your document. The color
  you defined will most likely not appear in the document. The reason for
  that is our way of handling framestyles. Since we can switch between
  framestyles with different background colors inside the document, we need
  a way to adapt the text colors as well. For this we defined 2 color themes
  which are automatically loaded when switching framestyles.

  \vspace{3mm}
  \begin{tabular}{ll}
    TUM default & color setup for white pages                  \\
    TUM inverse & color setup for blue, shaded, and flag pages \\
  \end{tabular}
\end{frame}

\begin{frame}[fragile]{Color themes}
  If you want to change the color of a particular element, you have to add
  the definition to one of the color themes. For this we offer the macro
  \lstinline|\addtoTUMcolortheme{<name>}{<code>}|.

  \vspace{3mm}
  As an example we can chage the title color on blue background pages with the
  command
  \begin{lstlisting}
    \addtoTUMcolortheme{TUM inverse}{
      \setbeamercolor{title}{fg=green}
    }
  \end{lstlisting}
\end{frame}


\part{Test document}
\section{Text only page}
\begin{frame}{Text only page}{with a subtitle}
  This is a text only page, to see if the spacing is correct.
  \lipsum[1]
\end{frame}

\section{List environments}
\subsection{Itemize page}
\begin{frame}{Itemize test}
  \begin{itemize}
    \item This is the first point.
    \item This is a second point.
      \begin{itemize}
        \item This point has sub items.
        \item And there is also one sub item which is very long and
          hopefully spans two lines, so I can see if the alignment of
          the second line is correct.
        \item Here the last sub item.
      \end{itemize}
    \item After sub items there are more items.
    \item Here is another item with a very long text that should span at
      least two lines. So I am telling a bit Lorem Ipsum stuff, to get the
      lines full.
  \end{itemize}
\end{frame}

\subsection{Enumerate page}
\begin{frame}{Enumerate test}
  \begin{enumerate}
    \item This is the first point.
    \item This is a second point.
      \begin{enumerate}
        \item This point has sub items.
        \item And there is also one sub item which is very long and
          hopefully spans two lines, so I can see if the alignment of
          the second line is correct.
        \item Here the last sub item.
      \end{enumerate}
    \item After sub items there are more items.
    \item Here is another item with a very long text that should span at
      least two lines. So I am telling a bit Lorem Ipsum stuff, to get the
      lines full.
  \end{enumerate}
\end{frame}

\subsection{Description page}
\begin{frame}{Description test}
  \begin{description}
    \item[TUM] is an excellent university
    \item[EI] is an excellent faculty
    \item[Nano] is an excellent group
    \item[ICMST] is a conference on mostly scientific topics
    \item[Longer Label] is to see the alignment
    \item[Even longer label] is to see the alignment
  \end{description}
\end{frame}

\section{Figures and Tables}
\subsection{Figure page}
\begin{frame}{A nice figure}
  \begin{figure}
    \includegraphics[height=4cm]{resources/TUM_Flags_169}
    \caption{This is a picure of Flags. Of course the excellent TUM flag
      is in the front!}
    \label{fig:TUMflags}
  \end{figure}
\end{frame}

\subsection{Table page}
\begin{frame}{An important table}
  \begin{table}
    \caption{This table contains very important data.}
    \label{tab:test}
    \begin{tabular}{lrc}
      \toprule
      a & b & c \\
      \midrule
      d & e & f \\
      g & h & i \\
      1 & 2 & 3 \\
      \bottomrule
    \end{tabular}
  \end{table}
\end{frame}


\part{Another part}
\section{Beamer Blocks}
\subsection{Alert Block}
\begin{frame}{ALERT!}
  \begin{alertblock}{Watch out!}
    This is how an alert block looks like.
    \begin{equation}
      a^2 + b^2 = c^2
    \end{equation}
  \end{alertblock}
  Let's also try an \alert{alerted} text and an \emph{emphasized} text.
\end{frame}

\begin{frame}{ALERT!}
  On this slied, we test \alert<2>{alerts with overlays} in both text and
  math mode. The variables \alert<2>{$x$} and $\alert<2>{y}$ will be alerted
  by color change only and not by boldfacing, because this would change the
  meaning of the variable.

  Alerting in math mode should also work in math environments
  \begin{equation}
    \alert<2>{\nabla}\times\vec{B} = \mu_{0} \left(
      \vec{J} + \varepsilon_{0}\frac{\partial \vec{E}}{\alert<2>{\partial t}}
    \right),
  \end{equation}
  and when alerting text with \alert<2>{inline math $\sin(\alpha)$}.
\end{frame}

\subsection{Example and Standard Blocks}
\begin{frame}{Example and Standard Blocks}
  \begin{exampleblock}{The only example}
    Here is an example.
  \end{exampleblock}
  \begin{block}{This is a standard block}
    Here we see the content of a standard block.
  \end{block}
\end{frame}

\section{Columns}
\begin{frame}{Columns}
  \begin{columns}[t, onlytextwidth]
    \column{.49\textwidth}
    Let's split the text of this frame in two columns! This could look
    really good.
    \begin{itemize}
      \item We can include items
      \item another one
      \item and the last one
    \end{itemize}
    The second column should contain a figure.
    \column{.49\textwidth}
    The top alignment is checked.
    \begin{figure}
      \theTUMLogo[height=2cm]
      \caption{This is the TUM logo, obviously.}
    \end{figure}
  \end{columns}
  \vspace{3mm}
  And this is text after the column environment, which should span the
  full width again. And indeed it does! But vertical space must be included
  manually :/
\end{frame}
